\section{Appendix}
In the following, we show the results of additional experiments we performed regarding initialization methods, mutation operators, and hyperparameters. Moreover, we provide pseudocode for our random walk initialization and our calculation of splitting values for data properties.

\subsection{Comparison of Initialization Methods}

In Table~\ref{tab:init_comparison}, we compare our random walk initialization to three widely used initialization methods~\citep{Koza1992GP}: \textsc{Grow}, \textsc{Full}, and \textsc{RampedHalfHalf}. \textsc{Grow} generates trees where leaves can have different depths, \textsc{Full} generates trees where all leaves have the same depth, and \textsc{RampedHalfHalf} generates half of a population's individuals with the \textsc{Grow} or \textsc{Full} method. We observe that the performance of the three methods is similar, with \textsc{Full} having a slight edge, while our random walk initialization performs significantly better. We hypothesize that this can be attributed to the fact that our generated class expressions, i.e., combinations of atomic concepts and operators, do actually appear in the instance data while the other three methods generate random/artificial combinations unrelated to the instance data. Manual spot checks confirmed our hypothesis.

\subsection{Comparison of Mutation Operators}

Table~\ref{tab:mutation} compares the performance of different mutation operators which are part of the DEAP framework~\citep{Fortin2012DEAP}:  \texttt{mut\-Uniform}, \texttt{mut\-Shrink}, \texttt{mut\-Node\-Re\-place\-ment} and \texttt{mut\-Insert}.
\texttt{mut\-Uniform} uniformly randomly selects a node in the tree and replaces it by a randomly generated tree as described in Section~\ref{subsec:evolutionary-algorithm}; \texttt{mut\-Shrink} randomly selects two nodes in the tree which are on a path from the root to a leaf and replaces the subtree induced by the higher node by the subtree induced by the lower node. \texttt{mut\-Node\-Re\-place\-ment} randomly replaces a single node in the tree by a single random node; \texttt{mut\-Insert} randomly selects an inner node, replaces it with a random inner node, and adds the original subtree induced by the node as a child.

While the differences are small on most datasets, \texttt{mut\-Uniform} performs much better than other approaches on \textit{Hepatitis}. We attribute this to the fact that \texttt{mut\-Uniform} introduces new constructs, like universal restrictions, negations, and cardinality restrictions, which \texttt{mut\-Shrink} does not.
Similarly, \texttt{mut\-Node\-Re\-place\-ment} and \texttt{mut\-Insert} can only introduce them to a limited extent by replacing or adding single nodes. 
These constructs seem to be particularly important for the Hepatitis dataset where almost all the good solutions contain them. 
In case of \textit{Lymphography}, \texttt{mut\-Uniform} learns long class expressions with training scores exceeding $0.9$ (not shown in tables) which cannot be matched on test data (0.84 in Table~\ref{tab:mutation}), hinting at slight overfitting.

\begin{table}[tbp]
    \centering
	\caption{Different initialization methods in terms of $\boldsymbol{F_1}$-measure for 10-fold cross-validation.}%
	\label{tab:init_comparison}
	\setlength{\tabcolsep}{5.9pt}
	\footnotesize
	\begin{tabular}{@{}lcccc@{}}
			\toprule
			\textbf{Learning Problem} & \textbf{Rand.\ Walk} & \textbf{\textsc{Grow}} & \textbf{\textsc{Full}} & \textbf{\textsc{Ramped}} \\
			\midrule
			Carcinogenesis & \textbf{0.70} $\boldsymbol{\pm}$ \textbf{0.12} & $0.61 \pm 0.20$ & $0.63 \pm 0.19$ & $0.60 \pm 0.21$ \\
			
			Family & \textbf{1.00} $\boldsymbol{\pm}$ \textbf{0.01} & $0.86 \pm 0.14$ & $0.87 \pm 0.14$ & $0.87 \pm 0.13$\\
			
			Hepatitis & \textbf{0.79} $\boldsymbol{\pm}$ \textbf{0.08} & $0.63 \pm 0.20$ & $0.76 \pm 0.14$ & $0.67 \pm 0.15$\\
			
			Lymphography & \textbf{0.84} $\boldsymbol{\pm}$ \textbf{0.09} & $0.83 \pm 0.08$ & \textbf{0.84} $\boldsymbol{\pm}$ \textbf{0.09} &  $0.83 \pm 0.11$\\
			
			Mammographic & \textbf{0.81} $\boldsymbol{\pm}$ \textbf{0.06} & $0.78 \pm 0.08$ & $0.78 \pm 0.08$ & $0.78 \pm 0.08$\\
			 
			Mutagenesis & \textbf{1.00} $\boldsymbol{\pm}$ \textbf{0.00} & \textbf{1.00} $\boldsymbol{\pm}$ \textbf{0.00} & \textbf{1.00} $\boldsymbol{\pm}$ \textbf{0.00} & \textbf{1.00} $\boldsymbol{\pm}$ \textbf{0.00}\\
			
			NCTRER & \textbf{1.00} $\boldsymbol{\pm}$ \textbf{0.00} & \textbf{1.00} $\boldsymbol{\pm}$ \textbf{0.00} & \textbf{1.00} $\boldsymbol{\pm}$ \textbf{0.00} & \textbf{1.00} $\boldsymbol{\pm}$ \textbf{0.00}\\
			
			Premier League & \textbf{1.00} $\boldsymbol{\pm}$ \textbf{0.00} & 0.99 $\pm$ 0.04  & 0.98 $\pm$ 0.05 & 0.98 $\pm$ 0.04 \\
			
			Pyrimidine & \textbf{0.91} $\boldsymbol{\pm}$ \textbf{0.14} & $0.84 \pm 0.22$ & $0.84 \pm 0.22$ & $0.83 \pm 0.22$\\
			\bottomrule
	\end{tabular}
\end{table}
\begin{table}[tbp]
    \centering
	\caption{Different mutation operators in terms of $\boldsymbol{F_1}$-mea\-sure for 10-fold cross-validation.}%
	\label{tab:mutation}
	\setlength{\tabcolsep}{6pt}
	\footnotesize
	\begin{tabular}{@{}lcccc@{}}
			\toprule
			\textbf{Learning Problem} & \textbf{\texttt{mutUniform}} & \textbf{\texttt{mutShrink}} &
			\textbf{\texttt{mutN.Repl.}} & \textbf{\texttt{mutInsert}} \\
			\midrule
			Carcinogenesis & \textbf{0.70} $\boldsymbol{\pm}$ \textbf{0.12} & $0.65 \pm 0.17$& $0.68 \pm 0.15$& $0.66 \pm 0.14$ \\
			 
			Family & \textbf{1.00} $\boldsymbol{\pm}$ \textbf{0.01} & $0.96 \pm 0.06$ & $0.99 \pm 0.02$& $0.97 \pm 0.06$\\
			
			Hepatitis & \textbf{0.79} $\boldsymbol{\pm}$ \textbf{0.08} & $0.40 \pm 0.21$& $0.53 \pm 0.16$& $0.49 \pm 0.23$\\
			
			Lymphography & \textbf{0.84} $\boldsymbol{\pm}$ \textbf{0.09} & $0.81 \pm 0.09$ & $0.81 \pm 0.09$ & $0.83 \pm 0.09$\\
			
			Mammographic & \textbf{0.81} $\boldsymbol{\pm}$ \textbf{0.06} & $0.80 \pm 0.07$& $0.81 \pm 0.07$& $0.80 \pm 0.06$\\
			 
			Mutagenesis & \textbf{1.00} $\boldsymbol{\pm}$ \textbf{0.00} & $0.97 \pm 0.11$ & \textbf{1.00} $\boldsymbol{\pm}$ \textbf{0.00} & $0.97 \pm 0.11$ \\
			
			NCTRER & \textbf{1.00} $\boldsymbol{\pm}$ \textbf{0.00} & $0.99 \pm 0.02$ & \textbf{1.00} $\boldsymbol{\pm}$ \textbf{0.00} & $0.99 \pm 0.02$\\
			
			Premier League & \textbf{1.00} $\boldsymbol{\pm}$ \textbf{0.00}& $0.98 \pm 0.05$ & $0.99 \pm 0.03$ & $0.98 \pm 0.06$ \\
			
			Pyrimidine & 0.91 $\pm$ 0.14 & $0.91 \pm 0.15$& $0.92 \pm 0.13$ & \textbf{0.93} $\boldsymbol{\pm}$ \textbf{0.14}\\
			
			\bottomrule
	\end{tabular}
\end{table}
\begin{table}[tbp]
    \centering
	\caption{Different settings for $\boldsymbol{maxT}$ in terms of $\boldsymbol{F_1}$-mea\-sure for 10-fold cross-validation.}%
	\label{tab:maxt}
	\setlength{\tabcolsep}{6pt}
	\footnotesize
	\begin{tabular}{@{}lcccc@{}}
			\toprule
			\textbf{Learning Problem} & \textbf{maxT = 1} & \textbf{maxT = 2} &
			\textbf{maxT = 4} & \textbf{maxT = 6} \\
			\midrule
			Carcinogenesis  & $0.64 \pm 0.18$& \textbf{0.70} $\boldsymbol{\pm}$ \textbf{0.12} & $0.67 \pm 0.14$ & $0.68 \pm 0.14$  \\
			 
			Family & $0.93 \pm 0.08$ & 1.00 $\pm$ 0.01 & $0.98 \pm 0.05$ & \textbf{1.00} $\boldsymbol{\pm}$ \textbf{0.00}\\
			
			Hepatitis & $0.71 \pm 0.15$ & 0.79 $\pm$ 0.08 & $0.79 \pm 0.12$ & \textbf{0.82} $\boldsymbol{\pm}$ \textbf{0.10}\\
			
			Lymphography & \textbf{0.84} $\boldsymbol{\pm}$ \textbf{0.09} & \textbf{0.84} $\boldsymbol{\pm}$ \textbf{0.09} & \textbf{0.84} $\boldsymbol{\pm}$ \textbf{0.09} & \textbf{0.84} $\boldsymbol{\pm}$ \textbf{0.09}\\
			
			Mammographic & $0.79 \pm 0.08$ & $0.81 \pm 0.06$ & \textbf{0.82} $\boldsymbol{\pm}$ \textbf{0.05} & $0.81 \pm 0.07$\\
			 
			Mutagenesis & \textbf{1.00} $\boldsymbol{\pm}$ \textbf{0.00} & \textbf{1.00} $\boldsymbol{\pm}$ \textbf{0.00} & \textbf{1.00} $\boldsymbol{\pm}$ \textbf{0.00} & \textbf{1.00} $\boldsymbol{\pm}$ \textbf{0.00}\\
			
			NCTRER & \textbf{1.00} $\boldsymbol{\pm}$ \textbf{0.00} & \textbf{1.00} $\boldsymbol{\pm}$ \textbf{0.00} & \textbf{1.00} $\boldsymbol{\pm}$ \textbf{0.00} & \textbf{1.00} $\boldsymbol{\pm}$ \textbf{0.00}\\
			
			Premier League & \textbf{1.00} $\boldsymbol{\pm}$ \textbf{0.00} & \textbf{1.00} $\boldsymbol{\pm}$ \textbf{0.00} & $1.00 \pm 0.02$ & \textbf{1.00} $\boldsymbol{\pm}$ \textbf{0.00} \\
			
			Pyrimidine & $0.85 \pm 0.21$ & \textbf{0.91} $\boldsymbol{\pm}$ \textbf{0.14} & $0.91 \pm 0.15$ & $0.87 \pm 0.17$\\
			
			\bottomrule
	\end{tabular}
\end{table}

\subsection{Comparison of maxT Settings}

The hyperparameter $maxT$ considers the number of outgoing triples of positive examples (c.f., Section~\ref{subsec:initialization} and Section~\ref{subsec:pseudocode}). Table~\ref{tab:maxt} evaluates different values of the $maxT$ parameter. It shows that a value of two is needed since setting $maxT:=1$ yielded worse $F_1$-measures. However, increasing the parameter beyond $maxT:=2$ did not significantly increase the $F_1$-measure and resulted in longer solutions. 

\subsection{Pseudocode}
\label{subsec:pseudocode}

Algorithm~\ref{alg:initialization} depicts the main loop of our random walk initialization. 
The \texttt{PopulationFromExamples} method takes a set of positive examples $E^+$ to generate the initial population. Therefore, it randomly selects positive examples $e^+ \in E^+$ and calls \texttt{ConceptFromExample}, which performs a biased random walk on $e^+$ to generate a concept. 
Algorithm~\ref{alg:initialization-property} turns a role/object pair $(r, o)$ into an existential restriction in case of an object property. In case of a data property, it generates a Boolean value restriction or a min/max numeric value restriction. 
Algorithm~\ref{alg:data_properties} shows the pseudocode for the calculation of splitting values for data properties. If not otherwise specified, sampling in the algorithms is performed uniformly at random.

\clearpage
\begin{algorithm}
	\DontPrintSemicolon
	\caption{Generates the initial population by creating concepts from positives examples via biased random walks.}%
	\label{alg:initialization}
	\textbf{Input:} Knowledge base $\mathcal{K}$, Positive examples $E^{+}$, Max.\ adjacent triples $\mathit{maxT}$, Set of $splits$ for each data property\\
    	\textbf{Output:} Set of concepts $population$ \\
    	\SetKwFunction{FMain}{PopulationFromExamples}
    	\SetKwProg{Pn}{Function}{:}{\KwRet $population$}
    	\Pn{\FMain{$\mathcal{K}$, $E^+$, $maxT$, $splits$}}{%
    		Compute the frequencies $\mathit{ct}[C]$ of the types (=atomic concepts) $C$ of the positive examples $E^{+}$\;
    		$population = []$\;
    		\While{$len(population) <$ \textsc{PopulationSize}}{%
    		Uniformly randomly pick $e^{+} \in E^{+}$\;
    		$population.append(  $ \texttt{ConceptFromExample}${(\mathcal{K}, e^{+}, \mathit{ct}, \mathit{maxT}, \mathit{splits})})$\;		
    		}
    	}
    \BlankLine
    \BlankLine
	\textbf{Input:} Knowledge base $\mathcal{K}$, Positive example $e^+$, Type frequencies $\mathit{ct}$, Max.\ adjacent triples $\mathit{maxT}$, Set of $splits$ for each data property \\
	\textbf{Output:} Concept $conc$ \\
	\SetKwFunction{F}{ConceptFromExample}
	\SetKwProg{Fn}{Function}{:}{\KwRet $conc$}
	\Fn{\F{$\mathcal{K}$, $e^+$, $\mathit{ct}$, $\mathit{maxT}$, $\mathit{splits}$}}{%
		Sample an atomic concept $C$ from $\{C \ |\ \mathcal{K} \models C(e^+)\}$ weighted by the frequencies $ct$ \tcp*[r]{Step (\ref{item1-rw})}
		$conc = C$\;
		$R = \{r \ | \ \mathcal{K} \models r(e^+, \cdot)\}$ \tcp*[r]{Step (\ref{item2-rw})}
		Sample up to $maxT$ $roles$ from $R$
		without replacement\;
		$role\_obj = \emptyset$\;
		\For{\textbf{each} $r \in roles$}{%
			Sample $o$ from $\{o \ |\ \mathcal{K} \models r(e^+, o) \}$
			\;
			$role\_obj = role\_obj \cup \{(r, o) \}$\;
		}
		\If{$|role\_obj|<maxT$}{%
		    $num = maxT - |role\_obj|$ \;
			Sample up to $num$ $new\_role\_obj$ from $\{(r,o) \ | \ \mathcal{K} \models r(e^+, o) \wedge r \in R  \}$
			without replacement\;
			$role\_obj = role\_obj \cup new\_role\_obj$\;
		}
		\For($\qquad\qquad \; \; \; \ $\tcp*[h]{Step (\ref{item3-rw})}){\textbf{each} $(r, o) \in role\_obj$}{%
			Sample an operator $op$ from $\{\sqcup, \sqcap \}$
			\;
			\eIf{ $type(r) = DataProperty $ \textbf{or} $\{s \ | \ \mathcal{K} \models s(o, \cdot)\} = \emptyset $ \textbf{or} $random() < 0.5$ }{%
				    $conc = conc$ op \texttt{RoleObjToConc}($\mathcal{K}$, $(r, o)$, $splits$)\; 
			    } 
			 {%
			    Sample a role $s$ from $\{s \ | \ \mathcal{K} \models s(o, \cdot)\}$
			    \;
				Sample $v$ from $\{(v \ | \ \mathcal{K} \models s(o, v) \wedge v \neq e^+\}$
				\;
				$conc = conc \ op \ \exists r$.\texttt{RoleObjToConc}($\mathcal{K}$, $(s, v)$, $splits$)\; 
			 }
		}
	}
\end{algorithm}
\begin{algorithm}
	\DontPrintSemicolon
	\caption{Turns the role/object pair into a concept.}%
	\label{alg:initialization-property}
    \textbf{Input:} Knowledge base $\mathcal{K}$, Role/Object Pair $(r, o)$, Set of $splits$ for each data property\\
    	\textbf{Output:} Concept $conc$ \\
    	\SetKwFunction{FMain}{RoleObjToConc}
    	\SetKwProg{Pn}{Function}{:}{\KwRet $conc$}
    	\Pn{\FMain{$\mathcal{K}$, $(r, o)$, $splits$}}{%
    	
    	\uIf{$type(o) = Individual$ }{%
    	        Sample an atomic concept $D$ from $\{D \ | \ \mathcal{K} \models D(o) \}$, if the set is empty take Thing\;
				$conc = (\exists r.D)$\; 
			    } 
			 \uElseIf{$type(o) = Boolean$}{%
			 	$conc = (r = o)$\;   
			 }
			\uElseIf{$type(o) = numeric$}{%
			        Find the split $v \in splits[r]$ that is closest to o\;
			        \eIf{$o \geq v$ } {%
			            $conc = (r \geq v)$\;   
			        }{%
			            $conc = (r \leq v)$\;  
			        }
			 }
    	}
\end{algorithm}
\newpage
\begin{algorithm}
	\DontPrintSemicolon
	\caption{Calculation of splitting values for data properties.}%
	\label{alg:data_properties}

	\textbf{Input:} Knowledge base $\mathcal{K}$, Data properties $P$, Number of splits $k$, Positive examples $E^{+}$, Negative examples $E^{-}$ \\
	\textbf{Output:} Set of $splits$ for each data property $d \in P$ \\
	\SetKwFunction{FMain}{CalculateSplits}
	\SetKwProg{Pn}{Function}{:}{\KwRet $splits$}
	\Pn{\FMain{$\mathcal{K}$, $P$, $k$, $E^+$, $E^-$}}{%
		For each $d \in P$ set $splits[d] = \emptyset$\;
		$current\_sets = \{E^+ \cup E^-\}$\;
		\While{$ P \neq \emptyset $ \textbf{and} $current\_sets \neq \emptyset$}{%
			$next\_level = \emptyset$\;
			\For{\textbf{each} $d \in P$}{%
					\For{\textbf{each} $E \in current\_sets$}{%
						$V^E_d = \{v \ | \ \mathcal{K} \models d(e,v) \wedge e \in E \}$\;
						Sort $V^E_d$ in ascending order\;
						$\bar{V}^E_d = \{(v_i  + v_{i+1}) /2 \ | \ i \in \{0,1,\ldots,|V^E_d| - 2\}\}$\;
						Find the best split $\bar{v} \in \bar{V}^E_d$ on $E$ depending on the information gain\;
						$splits[d] = splits[d] \cup \{ \bar{v} \}$\;
						Let $E^L$ and $E^R$ be the sets produced by splitting $E$ on $\bar{v}$\;
						$next\_level = next\_level \cup \{E^L, E^R\}$\;
						\If{$|splits[d]| \geq k$}{%
						    $P.remove(d)$\;
						    $break$\;
					    }
				}
			}
			$current\_sets = next\_level$\;
			Sort $current\_sets$ in descending order on the remaining entropy of the sets and remove sets with $0$ entropy\;
		}
	}
\end{algorithm}